% !TEX root = ../Projektdokumentation.tex
\section{Projektplanung}
\label{sec:Projektplanung}


\subsection{Projektphasen}
\label{sec:Projektphasen}

Für die Umsetzung des Projekts wurden gemäß Projektantrag insgesamt 40
Stunden angesetzt, wovon maximal acht Stunden auf die Dokumentation
entfallen.
Die einzelnen Arbeitspakete umfassen jeweils nicht mehr als sechs
Stunden.
Der Durchführungszeitraum erstreckt sich von der Genehmigung durch den
IHK-Prüfungsausschuss bis zum 30.04.2026.

Aus der Tabelle \ref{tab:Zeitplanung} kann entnommen werden, in welche
Hauptphasen das Projekt gegliedert wurde.
Eine detailliertere Zeitplanung mit allen einzelnen Arbeitspaketen
findet sich im \Anhang{app:Zeitplanung}.
\tabelle{Zeitplanung (geplant)}{tab:Zeitplanung}{ZeitplanungKurz}

% TODO: nach Abschluss des Projekts ggfs. Abweichungen zwischen geplanter
% und tatsächlicher Zeit hier kurz erläutern (Details im Fazit).


\subsection{Ressourcenplanung}
\label{sec:Ressourcenplanung}

Für die Durchführung des Projekts werden folgende Ressourcen verwendet.

\subsubsection{Hardware}
Für den \ac{PBS} steht ein dedizierter Server mit mindestens 32\,GB
\ac{RAM} und ausreichend Speicherkapazität zur Verfügung.
Die Storage-Kapazität wird in der Analysephase auf Basis der
bestehenden \ac{VM}- und Container-Landschaft dimensioniert.
% TODO: konkrete Hardware (Modell, CPU, RAM, Disks) nach Auswahl ergänzen.

Für die Administration und Entwicklung wird der reguläre
Arbeitslaptop des Prüflings genutzt.

\subsubsection{Software}
Auf dem Backup-Server kommt der \ac{PBS} (aktuelle Version) auf Basis von
Debian Linux zum Einsatz.
Für den Datastore wird \texttt{ext4} auf dem Hardware-\ac{RAID}\,5
verwendet.
Die Anbindung an die bestehende \ac{PVE}-Umgebung erfolgt über die
\ac{PBS}-eigenen Mechanismen (Storage-Plugin, Authentifizierung über
das lokale \texttt{root}-Passwort des \ac{PBS} und TLS-Fingerprint).
% TODO: ggfs. konkrete Versionen ergänzen.

\subsubsection{Netzwerk}
Für die Anbindung des \ac{PBS} ist ein dediziertes Backup-\ac{VLAN}
vorgesehen, das in Abstimmung mit dem Netzwerkteam eingerichtet wird.
Erforderliche Firewall-Freischaltungen werden im Rahmen der
Beschaffungs- und Vorbereitungsphase abgestimmt.

\subsubsection{Personal}
Folgende Personen unterstützen das Projekt:
\begin{itemize}
    \item Frank Thommen (\ac{ODCF}) als fachlicher Betreuer und Auftraggeber
    \item das \ac{ODCF}-Administrationsteam als Übergabeempfänger und Sparringspartner für die Anbindung an die \ac{PVE}-Umgebung
    \item das Netzwerkteam des \ac{DKFZ} für die Einrichtung des Backup-\ac{VLAN}
\end{itemize}
Die Umsetzung selbst erfolgt durch den Prüfling.

\subsection{Entwicklungsprozess}
\label{sec:Entwicklungsprozess}

Das Projekt wird in einem iterativen Vorgehen umgesetzt: Nach jeder
abgeschlossenen Phase erfolgt eine kurze Abstimmung mit dem Betreuer,
um das weitere Vorgehen festzulegen und Anpassungen an der Planung
vorzunehmen.
Da es sich um ein Infrastrukturprojekt handelt, ist eine klare
Trennung zwischen Analyse, Konzeption, Umsetzung und Validierung
sinnvoll. Innerhalb der Umsetzung wird jedoch testgetrieben
gearbeitet, indem einzelne Konfigurationsschritte direkt durch
Backup- und Restore-Tests verifiziert werden.
